\documentclass{situationist-preprint}

\title{Smartphones and Extended Cognition}
\msid{MND-09-b}
\msfamily{MND-09}
\msdiscipline{Philosophy of Mind}
\mstype{Consequence}
\mskeywords{extended mind; smartphones; reliable coupling; endorsement; cognitive offloading}

\begin{document}
\maketitle

\begin{abstract}
MND-09 argues that cognitive processes extend into external resources reliably coupled to the agent, and that a personal phone used as Otto uses his notebook is part of its owner's memory. We derive the consequences for the smartphone in detail. A phone is not part of its owner's mind as a whole; particular stores on it are parts of particular memories, according as they meet the four conditions of reliable coupling. Notes, contacts and calendar entries that the owner has made or adopted typically qualify; search results, feeds and recommendations typically do not. We then derive two further consequences: interference with a coupled store is interference with the owner's memory, and the owner's beliefs can change when such a store is altered without her knowledge.
\end{abstract}

\section{Store by store}

MND-09-a states the conditions under which an external resource is reliably coupled to an agent as a store of information of a given kind: it is reliably available and typically invoked when relevant, its information is endorsed without further scrutiny, it is easily accessible, and its information is there because the agent put it there or adopted it (Definition~2.3, conditions (a)--(d)). Theorem~3.2 holds that dispositional beliefs and the processes that retrieve them extend into resources so coupled; Corollary~3.3 applies this to a personal phone used as Otto uses his notebook \citep{ClarkChalmers1998}.

The conditions are relative to a kind of information. A phone is many stores, and the theorem applies to each separately.

\begin{proposition}\label{p:stores}
Of the stores on a typical personal phone, those whose contents the owner entered or deliberately adopted, and on which she relies without checking, are parts of her memory; those whose contents she has not endorsed are not.
\end{proposition}

\begin{proof}
A phone carried on the person and consulted when the relevant information is needed meets conditions (a) and (c) for every store on it. Conditions (b) and (d) vary by store. A contacts list, a calendar or a notes file is filled by the owner or by her adoption of entries offered to her, and its contents are acted on without scrutiny, as remembered information is; it meets all four conditions, and Theorem~3.2 of MND-09-a applies. A search engine's results, a news feed and a recommendation list are filled by others; the owner scrutinizes them, or would on reflection, and has not endorsed them. They fail condition (d), and by Proposition~4.1 of MND-09-a they are not part of her memory, however reliably available.
\end{proof}

Intermediate cases are decided by the same conditions. A navigation application's map was not made by the owner, and fails (d); a list of saved places she has marked on it meets it. A photograph library she took is a store of her experiences that she has endorsed by taking and keeping the photographs; whether it functions as memory depends on whether she relies on it without scrutiny, which varies between owners. The point of the proposition is that the question whether a phone is part of the mind has no general answer: the answer is given store by store and owner by owner.

Empirical work on cognitive offloading supports the pattern. People who expect information to remain available on a device recall it less well and recall better where it is stored \citep{Sparrow2011}: the device and the agent's biological memory divide the work, as the parity principle predicts that functionally integrated stores will.

\section{Two further consequences}

\begin{corollary}\label{c:harm}
Unauthorized deletion or alteration of a coupled store on a phone is interference with its owner's memory, not merely with her property.
\end{corollary}

\begin{proof}
By Proposition~\ref{p:stores}, a coupled store is part of the owner's memory. Deleting its contents destroys dispositional beliefs that she has; altering them changes what she believes. That this is done to a part of her memory located outside her skull does not change what is done, by the principle of locational neutrality on which MND-09-a relies (Principle~P2).
\end{proof}

Corollary~\ref{c:harm} has been drawn as an ethical consequence of extended cognition \citep{CarterPalermos2016}. Its force here is descriptive: the classification of an act of interference follows from the metaphysics of memory before any question of its wrongness arises.

\begin{corollary}\label{c:alter}
If a coupled store is altered without the owner's knowledge, she can come to have beliefs that she never formed.
\end{corollary}

\begin{proof}
Her dispositional beliefs are partly constituted by the contents of the store. If the contents change, so do the beliefs. Since condition (b) means that she endorses retrieved information without scrutiny, nothing in her use of the store corrects the change.
\end{proof}

\begin{objection}
Corollary~\ref{c:alter} is a reductio. Beliefs cannot be implanted by editing a file.
\end{objection}

\begin{reply}
Beliefs can be implanted by editing biological memory, as the literature on false memory shows, and the corollary says no more than that the same holds of an extended store. The past-endorsement condition (d) is satisfied at the time of coupling, not at each change of content; after an alteration the store may cease to meet it, and an owner who discovers the alteration will withdraw her reliance on it. Until then, the altered store plays the role of her memory, and what it says is what she believes, just as Otto believes what his notebook says.
\end{reply}

\section{Conclusion}

A smartphone extends its owner's mind store by store: where its contents are the owner's own and relied on without checking, they are part of her memory, and where they are supplied by others, they are not. Interference with a coupled store is interference with memory, and alteration of one can change what its owner believes. These are consequences of the extended mind thesis that MND-09 defends, applied to the device that most often meets its conditions.

\begin{thebibliography}{9}\small
\bibitem[Carter and Palermos(2016)]{CarterPalermos2016} Carter, J.~A. and Palermos, S.~O. (2016). Is having your computer compromised a personal assault? The ethics of extended cognition. \emph{Journal of the American Philosophical Association}, 2(4), 542--560.
\bibitem[Clark and Chalmers(1998)]{ClarkChalmers1998} Clark, A. and Chalmers, D.~J. (1998). The extended mind. \emph{Analysis}, 58(1), 7--19.
\bibitem[Sparrow et~al.(2011)]{Sparrow2011} Sparrow, B., Liu, J. and Wegner, D.~M. (2011). Google effects on memory: cognitive consequences of having information at our fingertips. \emph{Science}, 333(6043), 776--778.
\end{thebibliography}

\end{document}
